\startSection
	\selectlanguage{romanian}
	\sectionTitle{Imnul Național al României}
	
	\begin{question}{Imnul național al României, cunoscut sub titlul \enquote{Deșteaptă-te, române}, are o istorie profund legată de mișcările revoluționare din secolul al XIX-lea. Cine este autorul versurilor acestui imn, care exprimă chemarea la unitate și libertate?}
		\choice{Vasile Alecsandri}
		\choice{Mihai Eminescu}
		\choice[correct]{Andrei Mureșanu}
		\choice{George Coșbuc}
		\choice{Ion Heliade Rădulescu}
	\end{question}

	
	\begin{question}{Melodia imnului \enquote{Deșteaptă-te, române} a fost compusă pentru a însoți versurile lui Andrei Mureșanu. Cine este compozitorul care a creat această muzică devenită simbol național?}
		\choice{Ciprian Porumbescu}
		\choice[correct]{Anton Pann}
		\choice{Tiberiu Brediceanu}
		\choice{George Enescu}
		\choice{Dinu Lipatti}
	\end{question}

	
	\begin{question}{Imnul actual al României a fost adoptat oficial într-un moment de schimbare politică majoră. În ce an a devenit \enquote{Deșteaptă-te, române} imnul național al țării?}
		\choice{1918}
		\choice{1947}
		\choice{1989}
		\choice[correct]{1990}
		\choice{2001}
	\end{question}

	
	\begin{question}{Versurile imnului național românesc exprimă o chemare la trezirea conștiinței naționale. Care este primul vers cu care începe acest imn?}
		\choice{\enquote{Români, uniți-vă în cuget și-n simțiri}}
		\choice{\enquote{Sculați-vă din somnul cel de moarte}}
		\choice[correct]{\enquote{Deșteaptă-te, române, din somnul cel de moarte}}
		\choice{\enquote{Trăiască România liberă și unită}}
		\choice{\enquote{Ridică-te, nație mândră și curajoasă}}
	\end{question}

	
	\begin{question}{Înainte de adoptarea imnului \enquote{Deșteaptă-te, române}, România a avut mai multe imnuri oficiale. Care dintre următoarele NU a fost niciodată imn național al României?}
		\choice{\enquote{Trăiască Regele}}
		\choice{\enquote{Pe-al nostru steag e scris Unire}}
		\choice{\enquote{Trei culori}}
		\choice[correct]{\enquote{Hora Unirii}}
		\choice{\enquote{Deșteaptă-te, române}}
	\end{question}
\renderSection
